% ============================================================================
% COURS 10 — ÉLECTRONIQUE DE PUISSANCE
% Reprise progressive à partir de 14-ElectroniquePuissance.docx
% ============================================================================

\begin{savoirs}
\begin{itemize}
\item Les différentes natures de sources, courant et tension, ainsi que leurs règles d'association.
\item Les caractéristiques d'un interrupteur à l'état passant, tension nulle, et à l'état bloqué, courant nul.
\item La forme du courant de la charge d'un convertisseur statique pour les charges à courant constant, $L$--$E$ et $R$--$L$--$E$.
\item La définition du rapport cyclique.
\item Les définitions des valeurs moyenne et efficace.
\item Les principaux composants utilisés en électronique de puissance : diode, MOSFET et IGBT.
\item L'association transistor--diode antiparallèle permettant d'obtenir un interrupteur réversible en courant.
\item Les schémas de principe des hacheurs série, deux quadrants et quatre quadrants.
\item La démarche permettant de tracer les tensions et les courants d'un convertisseur.
\item La démarche permettant de déterminer la relation entre la tension d'entrée et la tension de sortie par les valeurs moyennes.
\item La démarche permettant de déterminer l'ondulation du courant et la valeur du rapport cyclique pour laquelle elle est maximale.
\end{itemize}
\end{savoirs}

\begin{savoirfaire}
\begin{itemize}
\item Tracer l'allure du courant d'une charge à courant constant, $L$--$E$ ou $R$--$L$--$E$.
\item Déterminer le rapport cyclique à partir d'un chronogramme.
\item Calculer une valeur moyenne ou efficace à partir d'un chronogramme.
\item Reconnaître une diode, un transistor MOSFET ou un IGBT sur un schéma électrique.
\item Tracer les tensions et les courants d'un convertisseur DC--DC.
\item Déterminer la relation entre les tensions d'entrée et de sortie d'un convertisseur DC--DC.
\item Déterminer l'expression de l'ondulation de courant.
\item Identifier les composants passants à partir des chronogrammes de tension et de courant dans la charge.
\end{itemize}
\end{savoirfaire}

\section{Généralités sur l'électronique de puissance}

\subsection{Introduction}

L'énergie électrique se présente sous deux formes :
\begin{itemize}
\item continue, ou DC, comme pour les piles et les batteries ;
\item alternative, ou AC, comme pour les réseaux de distribution.
\end{itemize}

On parle d'\textbf{électronique de puissance} lorsqu'on étudie les \textbf{convertisseurs statiques d'énergie} permettant de modifier la présentation de l'énergie selon les besoins : chargeur de batterie, interconnexion de réseaux, variation de vitesse d'un moteur, alimentation électronique, etc.

On distingue ainsi :
\begin{itemize}
\item l'électronique des courants faibles, qui traite des informations codées sous forme de signaux électriques ;
\item l'électronique de puissance, ou courants forts, qui traite de la transformation et de la distribution de l'énergie électrique.
\end{itemize}

\begin{center}
\TLChaineConvertisseursPuissance
\end{center}

\subsubsection*{Exemples d'utilisation}

L'électronique de puissance est présente dans de nombreux systèmes, pour des puissances allant du milliwatt au mégawatt.

\paragraph{Alimentation d'un ordinateur.}
À partir du secteur $230\,\mathrm{V}$ alternatif, l'alimentation fournit plusieurs tensions continues destinées aux différents constituants électroniques.

\paragraph{TGV PBKA.}
Ce train doit pouvoir fonctionner sous plusieurs alimentations : $1500\,\mathrm{V}$ continu, $3000\,\mathrm{V}$ continu, $15\,\mathrm{kV}$ à $16\tfrac{2}{3}\,\mathrm{Hz}$ et $25\,\mathrm{kV}$ à $50\,\mathrm{Hz}$.

\paragraph{Variation de vitesse d'une MCC.}
Une solution résistive simple permettrait de faire varier la tension appliquée au moteur, mais elle présente deux défauts majeurs :
\begin{itemize}
\item la tension de sortie dépend de la charge ;
\item le rendement est très faible, car une part importante de l'énergie est dissipée par effet Joule.
\end{itemize}

Pour obtenir un rendement élevé, les convertisseurs statiques utilisent des interrupteurs électroniques, commandés alternativement dans leurs états passant et bloqué.

\subsection{Composants de base d'un convertisseur statique : bras de pont}

Le composant élémentaire d'un convertisseur statique est l'\textbf{interrupteur}. Dans le modèle idéal :
\[
\begin{array}{c|cc}
 & u_K & i_K \\
\hline
\text{passant} & 0 & \text{imposé par le circuit} \\
\text{bloqué} & \text{imposé par le circuit} & 0
\end{array}
\]

\begin{center}
\TLBrasDePont
\end{center}

L'interrupteur idéal ne dissipe donc aucune puissance dans ses deux états stables :
\[
p_K=u_K i_K=0.
\]
Les pertes réelles apparaissent pendant les conductions non idéales et surtout pendant les commutations.

\subsection{Nature des sources}

\paragraph{Source de tension.}
Une source de tension impose sa tension ; le courant dépend du circuit extérieur. Elle ne doit jamais être court-circuitée.

\begin{center}
\begin{circuitikz}[european resistors]
\draw (0,0) node[eground]{} to[V,l=$E$] (0,3) -- (2.2,3);
\draw (2.2,0) -- (0,0);
\draw[->] (2.2,2.7) -- (2.2,.3) node[midway,right]{$u=E$};
\end{circuitikz}
\end{center}

\paragraph{Source de courant.}
Une source de courant impose son courant ; la tension dépend du circuit extérieur. Elle ne doit jamais être ouverte brutalement.

\begin{center}
\begin{circuitikz}
\draw (0,0) node[eground]{} to[I,l=$I$] (0,3) -- (2.2,3);
\draw (2.2,0) -- (0,0);
\draw[->] (1.1,3.25) -- (2.1,3.25) node[midway,above]{$i=I$};
\end{circuitikz}
\end{center}

\subsection{Règles d'association des sources}

Les associations interdites sont :
\begin{itemize}
\item la mise en parallèle de deux sources de tension idéales de valeurs différentes ;
\item la mise en série de deux sources de courant idéales de valeurs différentes ;
\item le court-circuit d'une source de tension ;
\item l'ouverture d'une source de courant.
\end{itemize}

Dans un convertisseur statique, les commutations doivent donc toujours respecter la nature instantanée des sources placées de part et d'autre du convertisseur.

\subsection{Modification de la nature apparente d'une source}

Une inductance s'oppose aux variations brusques du courant :
\[
u_L(t)=L\frac{\mathrm d i_L}{\mathrm dt}.
\]
Le courant dans une inductance est continu. Placée en série avec une source de tension, elle donne à l'ensemble un comportement de source de courant sur la durée d'une commutation.

\begin{center}
\begin{circuitikz}[european resistors]
\draw (0,0) node[eground]{} to[V,l=$E$] (0,3)
      to[L,l=$L$,i>^=$i_L$] (3.5,3)
      to[short,-o] (4.5,3);
\draw (4.5,0) node[ocirc]{} -- (0,0);
\draw[<->] (4.7,0)--node[right]{$u_s$}(4.7,3);
\end{circuitikz}
\end{center}

Réciproquement, un condensateur s'oppose aux variations brusques de tension :
\[
i_C(t)=C\frac{\mathrm d u_C}{\mathrm dt}.
\]
La tension à ses bornes est continue. Placé en parallèle, il donne à l'ensemble un comportement de source de tension pendant la commutation.

\subsection{Forme du courant pour différents modèles de charge}

\paragraph{Charge à courant constant.}
Le courant est supposé parfaitement constant. Ce modèle constitue l'approximation limite d'une charge très inductive.

\paragraph{Charge $L$--$E$.}
Une inductance $L$ est associée à une force électromotrice $E$. Ce modèle est adapté à une machine à courant continu lorsque la résistance d'induit est négligée.

\begin{center}
\begin{circuitikz}[european resistors]
\draw (0,2.5) node[left]{$u_s$} to[short,o-] (1,2.5)
      to[L,l=$L$,i>^=$i$] (3.5,2.5)
      to[V,l=$E$] (3.5,0) -- (0,0) node[ocirc]{};
\end{circuitikz}
\end{center}

La loi de la maille donne :
\[
u_s=L\frac{\mathrm di}{\mathrm dt}+E.
\]
Le courant évolue donc linéairement lorsque $u_s$ est constant.

\paragraph{Charge $R$--$L$--$E$.}
Le modèle complet d'une MCC ajoute la résistance d'induit :

\begin{center}
\begin{circuitikz}[european resistors]
\draw (0,2.5) node[left]{$u_s$} to[short,o-] (1,2.5)
      to[R,l=$R$] (2.8,2.5)
      to[L,l=$L$,i>^=$i$] (5,2.5)
      to[V,l=$E$] (5,0) -- (0,0) node[ocirc]{};
\end{circuitikz}
\end{center}

La loi de la maille s'écrit :
\[
u_s=Ri+L\frac{\mathrm di}{\mathrm dt}+E.
\]
Le courant n'est plus affine par morceaux, mais suit une évolution exponentielle autour de sa valeur d'équilibre.

\subsection{Commutations possibles}

Une commutation est dite \textbf{spontanée} lorsqu'elle est imposée par l'évolution des grandeurs électriques. Elle est dite \textbf{commandée} lorsqu'un signal de commande impose le changement d'état du composant.

Dans un bras de pont, les deux interrupteurs ne doivent jamais être passants simultanément, sous peine de court-circuiter la source de tension. Ils ne doivent pas non plus être simultanément bloqués lorsque la charge se comporte comme une source de courant.

\subsection{Rapport cyclique et découpage}

Pour un signal périodique de période $T$, passant à l'état haut pendant $\alpha T$, le rapport cyclique est :
\[
\alpha=\frac{t_{\mathrm{on}}}{T},\qquad 0\leq\alpha\leq1.
\]

\begin{center}
\TLChronogrammePWM
\end{center}

\subsection{Valeur moyenne}

Pour une grandeur périodique $x(t)$ de période $T$ :
\[
X_{\mathrm{moy}}=\frac{1}{T}\int_0^T x(t)\,\mathrm dt.
\]
La valeur moyenne caractérise la composante continue du signal.

\subsection{Valeur efficace}

La valeur efficace est définie par :
\[
X_{\mathrm{eff}}=\sqrt{\frac{1}{T}\int_0^T x^2(t)\,\mathrm dt}.
\]
Elle correspond à la valeur continue produisant le même effet thermique dans une résistance.

\section{Composants en électronique de puissance}
Les composants usuels sont la diode, le transistor MOSFET, l'IGBT et le thyristor. On distingue :
\begin{itemize}
\item la caractéristique statique, qui décrit les états passant et bloqué ;
\item la caractéristique dynamique, qui décrit les commutations spontanées ou commandées.
\end{itemize}

La diode commute spontanément selon les conditions électriques. Les transistors MOSFET et IGBT sont commandés à l'amorçage et au blocage. L'association d'un transistor et d'une diode antiparallèle donne un interrupteur réversible en courant.

\subsection{Pertes et dimensionnement}
Les pertes par conduction s'évaluent à partir de la chute de tension à l'état passant et du courant. Les pertes de commutation dépendent de la fréquence de découpage et des énergies dissipées à l'amorçage et au blocage :
\[
P_{\mathrm{com}}=f_s\left(E_{\mathrm{on}}+E_{\mathrm{off}}\right).
\]
Le dimensionnement impose de vérifier les contraintes maximales en tension, courant, puissance dissipée et température de jonction.

\section{Méthode d'étude d'un convertisseur DC--DC}
La démarche est systématique :
\begin{enumerate}
\item identifier les séquences de fonctionnement sur une période ;
\item remplacer chaque interrupteur par son modèle passant ou bloqué ;
\item écrire les lois électriques dans chaque séquence ;
\item tracer tensions et courants ;
\item appliquer la valeur moyenne nulle de la tension aux bornes d'une inductance en régime périodique ;
\item calculer la valeur moyenne de sortie et l'ondulation du courant ;
\item déterminer les contraintes sur les composants.
\end{enumerate}

\section{Hacheur série ou Buck}
Le hacheur série convertit une tension continue en une tension moyenne plus faible.
\begin{center}\TLHacheurBuck\end{center}

Pour $0<t<\alpha T$, $K$ est passant et $u_L=E-u_s$. Pour $\alpha T<t<T$, $K$ est bloqué et la diode assure la continuité du courant : $u_L=-u_s$.
En régime périodique :
\[
\langle u_L\rangle=0
\quad\Longrightarrow\quad
\boxed{U_s=\alpha E}.
\]
L'ondulation crête-à-crête du courant vaut :
\[
\Delta I=\frac{(E-U_s)\alpha T}{L}
=\frac{E\alpha(1-\alpha)}{Lf_s}.
\]
Elle est maximale pour $\alpha=\tfrac12$.

\section{Hacheur parallèle ou Boost}
Le hacheur parallèle élève la tension moyenne de sortie.
\begin{center}\TLHacheurBoost\end{center}

Lorsque $K$ est passant, l'inductance emmagasine de l'énergie. Lorsqu'il est bloqué, l'inductance restitue cette énergie à la charge à travers la diode. En régime périodique idéal :
\[
\boxed{U_s=\frac{E}{1-\alpha}}.
\]

\section{Hacheur deux quadrants}
Le hacheur deux quadrants permet une tension de sortie de signe fixé et un courant réversible. Il convient notamment aux phases motrices et génératrices d'une MCC.
\begin{center}\TLHacheurDeuxQuadrants\end{center}

Les commandes de $K_1$ et $K_2$ sont complémentaires afin d'éviter le court-circuit de la source. Selon le signe du courant, la conduction est assurée par un transistor ou par sa diode antiparallèle.

\section{Hacheur quatre quadrants}
Le pont complet autorise les deux signes de tension et de courant. Il permet donc le fonctionnement dans les quatre quadrants du plan $(u_s,i_s)$.
\begin{center}\TLHacheurQuatreQuadrants\end{center}

Avec une commande bipolaire, la tension moyenne de sortie peut s'écrire :
\[
\boxed{U_s=(2\alpha-1)E}.
\]

\section{Conversion indirecte}
Les convertisseurs isolés utilisent un étage haute fréquence et un transformateur. Ils permettent l'adaptation du niveau de tension, l'isolement galvanique et parfois plusieurs sorties.

\section{Convertisseurs AC--DC}
Le redressement transforme une tension alternative en une tension unidirectionnelle. Le montage monophasé à pont de diodes, ou PD2, utilise quatre diodes.

\begin{center}\TLPontDiodes\end{center}

Pour une charge résistive idéale et une entrée sinusoïdale $u_e(t)=\widehat U\sin(\omega t)$ :
\[
u_s(t)=|u_e(t)|,
\qquad
U_{s,\mathrm{moy}}=\frac{2\widehat U}{\pi}.
\]

\section{Puissances en régime périodique}
La puissance instantanée est $p(t)=u(t)i(t)$. La puissance active est :
\[
P=\frac1T\int_0^T u(t)i(t)\,\mathrm dt.
\]
En régime sinusoïdal :
\[
P=UI\cos\varphi,\qquad Q=UI\sin\varphi,\qquad S=UI.
\]
En présence d'harmoniques, la relation générale devient :
\[
S^2=P^2+Q^2+D^2,
\]
où $D$ est la puissance déformante. Le facteur de puissance est :
\[
\lambda=\frac{P}{S}.
\]

\section{Transformateur monophasé parfait}
Pour un transformateur idéal :
\[
\frac{U_2}{U_1}=\frac{N_2}{N_1}=m,
\qquad
\frac{I_2}{I_1}=\frac1m,
\qquad
P_1=P_2.
\]

\section{Redresseur à absorption sinusoïdale}
L'objectif est de rendre le courant absorbé quasi sinusoïdal et en phase avec la tension du réseau afin d'améliorer le facteur de puissance. Une structure classique associe un pont redresseur puis un étage Boost commandé. Des structures mono-étage réalisent simultanément redressement et correction du facteur de puissance.

\section{Synthèse}
\begin{itemize}
\item Le respect de la nature des sources conditionne la structure du convertisseur.
\item L'étude repose sur les séquences de commutation et les modèles équivalents.
\item Le Buck abaisse la tension, le Boost l'élève.
\item Les hacheurs deux et quatre quadrants rendent l'énergie réversible.
\item Le pont PD2 réalise un redressement double alternance.
\item Le dimensionnement doit intégrer contraintes électriques, pertes et thermique.
\end{itemize}
